\section{Glossario dei Termini di Dominio}

La presente sezione definisce in modo non ambiguo il lessico specialistico e i concetti fondamentali impiegati nell'intero documento di specifica, al fine di garantire uniformità terminologica tra stakeholder, analisti e sviluppatori.

\begin{center}
\renewcommand{\arraystretch}{1.3}
\begin{longtable}{|p{4.2cm}|p{11.2cm}|}
\hline
\rowcolor{tableheader} \textbf{Termine} & \textbf{Definizione e Significato nel Sistema} \\ \hline
\endfirsthead

\hline
\rowcolor{tableheader} \textbf{Termine} & \textbf{Definizione e Significato nel Sistema} \\ \hline
\endhead

\textbf{MyAma} & Piattaforma digitale integrata per la gestione e prenotazione dei servizi di smaltimento di rifiuti ingombranti e speciali. \\ \hline

\textbf{Cittadino / Cliente} & Utente esterno principale del sistema che effettua richieste di smaltimento di rifiuti ingombranti. \\ \hline

\textbf{Autista AMA} & Dipendente operativo AMA incaricato dello svolgimento dei ritiri a domicilio mediante veicoli aziendali. \\ \hline

\textbf{Operatore di Sede} & Dipendente operativo AMA assegnato a un centro di raccolta (sede) per la gestione dei conferimenti diretti. \\ \hline

\textbf{Lavoratore AMA} & Categoria generale che raggruppa Autisti e Operatori di sede, caratterizzata da competenze, turni e mansioni. \\ \hline

\textbf{Ritiro a domicilio} & Servizio di prelievo del rifiuto ingombrante presso il domicilio indicato dal Cittadino tramite autista e mezzo aziendale. \\ \hline

\textbf{Conferimento in sede} & Consegna diretta del rifiuto ingombrante da parte del Cittadino presso una specifica sede / isola ecologica AMA. \\ \hline

\textbf{Prenotazione} & Entità informativa che formalizza la richiesta di un servizio (ritiro o conferimento), specificando rifiuto, slot temporale, luogo e attori. \\ \hline

\textbf{Stato Prenotazione} & Ciclo di vita della prenotazione (\textit{In attesa, Confermata, In corso, Completata, Annullata, Non eseguita}). \\ \hline

\textbf{Rifiuto Ingombrante} & Bene o materiale non smaltibile nei circuiti ordinari (es. mobili, grandi elettrodomestici, sanitari, RAEE). \\ \hline

\textbf{Tipologia di Rifiuto} & Categoria merceologica del rifiuto utile a stabilire vincoli di smaltimento, ingombro e mezzi idonei. \\ \hline

\textbf{Sede AMA / Isola Ecologica} & Struttura fisica territoriale preposta alla ricezione dei rifiuti conferiti dai cittadini. \\ \hline

\textbf{CAP / Zona Servita} & Suddivisione territoriale che determina la copertura del servizio e l'assegnazione delle sedi/autisti competenti. \\ \hline

\textbf{Veicolo / Mezzo} & Automezzo AMA utilizzato per i ritiri a domicilio, caratterizzato da portata massima in peso e volume. \\ \hline

\textbf{Disponibilità} & Slot temporale (data e fascia oraria) in cui il servizio è erogabile in base alla disponibilità di personale e mezzi. \\ \hline

\textbf{Assegnazione} & Associazione formale tra una prenotazione di ritiro e le risorse operative (Autista, Mezzo, Itinerario). \\ \hline

\textbf{Esito del Servizio} & Registrazione finale conclusiva dell'operazione di smaltimento (\textit{Completato con successo, Rifiutato/Non conforme, Assente}). \\ \hline
\end{longtable}
\end{center}
